% =============================================================================
% UNIVERSITY OF BRISTOL MSc DATA SCIENCE GROUP DISSERTATION
% Working master prepared directly from the supplied University template:
% DataScience_UoB_MSc_thesis(2)(2).tex
%
% IMPORTANT
% - Keep the University-supplied dissertation.cls unchanged.
% - Keep the University logo file as logo_uob_color.eps (or equivalent basename).
% - Keep the BibTeX workflow using sample_bibtex.bib unless/until the final
%   bibliography database is renamed deliberately.
% - The MSc dissertation-type field is intentionally left blank because the
%   required X/Y/Z project-type percentage has not yet been verified.
% =============================================================================

% The document class supplies options to control rendering of standard features.
% The University template specifies:
% - author names as FirstName LastName;
% - supervisor as Title FirstName LastName;
% - degree as BSc, MEng, MSci, MSc or PhD;
% - year as the four-digit year of submission.

\documentclass[
                    % the names of the group authors
                    author={\shortstack{Tanishk Nanasaheb Shinde\\Megh Kashilkar\\Sanskar Sanju Gade\\Jitendra Suwalka}},
                    % the name of the supervisor
                    supervisor={\begin{tabular}[t]{@{}l@{}}Dr. Felipe Campelo\\Assistant Supervisors: Dr. Ahmad Al-Mallahi\\(Dalhousie University / McCain Foods representative) and Jason McLaren\end{tabular}},
                    % the degree programme
                    degree={MSc},
                    % the dissertation title
                    title={Edge-Deployable Rock Detection and Burial Status Classification},
                    % optional subtitle
                    subtitle={},
                    % MSc dissertation type: replace with verified X/Y/Z value if required
                    type={},
                    % year of submission
                      year={2026}]{dissertation}

% Additional packages required by the report content. The University-supplied
% dissertation class remains unchanged.
\usepackage{booktabs}
\usepackage{array}
\usepackage{placeins}

% The report is supplied electronically and must not contain automatically
% inserted blank verso pages.
\let\cleardoublepage\clearpage

% The supplied class prints \today on the title page and declaration.
% Fix it to the documented submission deadline; edit this line if the actual
% submission date is different.
\renewcommand{\today}{11th September, 2026}

% -----------------------------------------------------------------------------
% GROUP DECLARATION ADAPTATION
% -----------------------------------------------------------------------------
% The University-supplied class was written for a single author and its
% declaration says "the work of the Author". The class file itself is left
% untouched. This local override changes only the authorship grammar so the
% group report does not present collective work as belonging to one person.
% Confirm the exact declaration wording with the programme team before final
% submission if a separate official group-declaration form is supplied.

\makeatletter
\renewcommand{\makedecl}{
  \thispagestyle{plain}

  \chapter*{Declaration}

  This dissertation is submitted to the University of Bristol in accordance
  with the requirements of the degree of {\dissertation@degree} in the Faculty
  of Engineering. It has not been submitted for any other degree or diploma
  of any examining body. Except where specifically acknowledged, it is all
  the work of the Authors.

  \vspace{5cm}

  \noindent
  \begin{tabular}{l}
  Tanishk Nanasaheb Shinde\\
  Megh Kashilkar\\
  Sanskar Sanju Gade\\
  Jitendra Suwalka
  \end{tabular}
  \hfill \today
}
\makeatother

\setlength{\emergencystretch}{2em}

% In-flow chart used for the supplied replacement visualisations. The optional
% first argument controls width so paired figures can remain close to the text.
\newcommand{\reportchart}[4][0.82\textwidth]{%
  \begin{figure}[!htbp]
  \centering
  \includegraphics[width=#1,keepaspectratio]{#2}
  \caption{#3}
  \label{#4}
  \end{figure}
}

\begin{document}

% =============================================================================
% TITLE PAGE
% =============================================================================

\hypersetup{pageanchor=false}
\maketitle

% =============================================================================
% FRONT MATTER
% =============================================================================

\frontmatter
\hypersetup{pageanchor=true}

\makedecl

\tableofcontents
\listoffigures
\listoftables
% No algorithms or source-code listings are currently planned for the main
% dissertation. Uncomment these if they are later used.
% \listofalgorithms
% \lstlistoflistings

% -----------------------------------------------------------------------------

\chapter*{Abstract}

\noindent
Agricultural rocks can interfere with field operations and potentially damage machinery,
motivating computer-vision systems capable of detecting their location and burial status.
This dissertation investigates the deployment of a two-class YOLOv8n detector for
identifying completely exposed and half-buried rocks in agricultural RGB imagery on a
Raspberry Pi 5. The work focuses on the transition from model development to
resource-constrained edge inference rather than implementing the wider GPS mapping or
automated rock-removal system.

The original dataset contained 448 images and 706 annotated rocks. YOLOv8n was selected
from YOLOv8n, YOLO11n and YOLO11s using validation mAP50--95 as the predefined primary
selection criterion. The selected detector was then prepared in six deployment
representations: PyTorch FP32, ONNX FP32, TFLite FP32, TFLite INT8, NCNN FP32 and NCNN
FP16. A common 90-image, 135-rock, $640 \times 640$ workload was used to evaluate
detection quality, inference latency, throughput, CPU utilisation, process memory,
rock-count error, temperature and operational behaviour on Raspberry Pi 5 CPU hardware.

Five representations completed the full benchmark. PyTorch FP32 achieved the strongest
aggregate edge-benchmark detection quality, with mAP50 of 0.548, but required 397.48~ms
mean latency and achieved only 2.52 median FPS. TFLite INT8 reduced mean latency to
44.13~ms and increased median throughput to 22.62~FPS, while mAP50 decreased to 0.503.
NCNN FP16 provided an alternative accuracy--efficiency compromise, achieving 13.82 median
FPS with mAP50 of 0.537. ONNX FP32 supported single-image inference but did not complete
sustained benchmarking.

TFLite INT8 additionally completed approximately 2,533 continuous inferences over two
minutes without recorded throttling or reboot, although temperature increased to
approximately 76--77~$^\circ$C. Half-buried rocks remained substantially more difficult
to recognise than completely exposed rocks, and this weakness existed before edge
conversion.

The results support the hypothesis that edge-optimised and reduced-precision
representations can substantially improve Raspberry Pi inference efficiency while
introducing runtime-dependent accuracy and operational trade-offs. TFLite INT8 was
therefore selected as the preferred representation for the tested CPU-only configuration.
The resulting system constitutes an academic edge-deployment proof-of-concept rather than
a production-ready agricultural solution.

% -----------------------------------------------------------------------------
% Summary of Changes is omitted because this is not a resubmission.
% -----------------------------------------------------------------------------

\chapter*{Supporting Technologies}

\noindent
The principal third-party hardware and software resources used in the project were:

\begin{itemize}
\setlength{\itemsep}{0.25em}
\item \textbf{Raspberry Pi 5 Model B} --- target edge-computing platform for CPU-only deployment and benchmarking of the selected YOLOv8n detector.
\item \textbf{Official Raspberry Pi 27~W USB-C power supply (5.1~V / 5~A)} --- used for the final Raspberry Pi benchmark configuration.
\item \textbf{Debian GNU/Linux 13} --- operating system used on the Raspberry Pi deployment platform.
\item \textbf{Python 3.13.5} --- principal Raspberry Pi programming environment for inference, benchmarking, verification and continuous-inference testing.
\item \textbf{Ultralytics 8.4.115} --- used for YOLO model handling, inference and preparation of alternative deployment representations.
\item \textbf{PyTorch} --- original FP32 representation and reference implementation for the selected detector.
\item \textbf{ONNX / ONNX Runtime 1.28.0} --- used for ONNX FP32 export and CPU inference; isolated inference succeeded but a valid sustained benchmark was not completed.
\item \textbf{TensorFlow Lite / LiteRT} --- used for FP32 and post-training INT8 deployment representations and on-device inference.
\item \textbf{NCNN} --- lightweight inference framework used for FP32 and FP16 deployment on the Raspberry Pi CPU.
\item \textbf{OpenCV 5.0.0} --- image loading, preprocessing and supporting computer-vision operations.
\item \textbf{NumPy 2.5.1} --- numerical and array-based processing in inference and evaluation scripts.
\item \textbf{Pandas 3.0.5} --- organisation and analysis of benchmark outputs and experimental results.
\item \textbf{psutil 7.2.2} --- CPU-utilisation and process-memory monitoring during Raspberry Pi benchmarking.
\item \textbf{CLIP ViT-B/32} --- separate crop-level burial-status classification research branch; it was not part of the verified Raspberry Pi benchmark.
\item \textbf{SSH and SCP} --- remote Raspberry Pi access and transfer of deployment files and benchmark outputs.
\end{itemize}

% -----------------------------------------------------------------------------

\chapter*{Notation and Acronyms}

\noindent
Important acronyms and notation used throughout the dissertation are summarised below.
Terms are also defined at first use in the main text.

\vspace{0.5cm}
\noindent
\textbf{Acronyms and technical abbreviations}

\vspace{0.3cm}
\begin{tabular}{p{2.7cm}p{11.1cm}}
AI & Artificial Intelligence \\
AP & Average Precision \\
CLIP & Contrastive Language--Image Pre-training \\
CPU & Central Processing Unit \\
F1 & Harmonic mean of precision and recall \\
FP16 & 16-bit floating-point numerical representation \\
FP32 & 32-bit floating-point numerical representation \\
FPS & Frames Per Second \\
GPS & Global Positioning System \\
INT8 & 8-bit integer numerical representation \\
IoU & Intersection over Union \\
MAE & Mean Absolute Error \\
mAP & mean Average Precision \\
mAP50 & mean Average Precision evaluated at IoU threshold 0.50 \\
mAP50--95 & mean Average Precision averaged across IoU thresholds from 0.50 to 0.95 in increments of 0.05 \\
NCNN & Lightweight neural-network inference framework for mobile and embedded deployment \\
ONNX & Open Neural Network Exchange \\
RAM & Random Access Memory \\
RGB & Red--Green--Blue image representation \\
RSS & Resident Set Size; process memory resident in physical RAM \\
SCP & Secure Copy Protocol \\
SPPF & Spatial Pyramid Pooling --- Fast \\
SSH & Secure Shell \\
TFLite & TensorFlow Lite \\
LiteRT & Google's on-device machine-learning runtime ecosystem used for TensorFlow Lite model execution \\
YOLO & You Only Look Once \\
\end{tabular}

\vspace{0.7cm}
\noindent
\textbf{Mathematical notation}

\vspace{0.3cm}
\begin{tabular}{p{2.7cm}p{11.1cm}}
$B_{\mathrm{pred}}$ & Predicted bounding box \\
$B_{\mathrm{gt}}$ & Ground-truth bounding box \\
$TP$ & Number of true-positive detections \\
$FP$ & Number of false-positive detections \\
$FN$ & Number of false-negative detections \\
$y_i$ & Ground-truth rock count for image $i$ \\
$\hat{y}_i$ & Predicted rock count for image $i$ \\
$N$ & Number of images included in an evaluation \\
$W,H$ & Width and height of the original image \\
$W',H'$ & Width and height after aspect-ratio-preserving resizing \\
$s$ & Letterbox image-scaling factor \\
$r$ & Real-valued quantity before quantisation \\
$q$ & Quantised integer representation \\
$S$ & Quantisation scale factor \\
$Z$ & Quantisation zero point \\
\end{tabular}

% -----------------------------------------------------------------------------

\chapter*{Acknowledgements}

\noindent
The authors gratefully acknowledge Dr. Felipe Campelo for his academic supervision,
regular guidance and constructive feedback throughout the project. His support helped the
group refine the research question, maintain a clear distinction between completed evidence
and planned work, and present the detector-development and edge-deployment results with the
appropriate level of critical evaluation.

The authors also thank the Assistant Supervisors, Dr. Ahmad Al-Mallahi and Jason McLaren,
for helping to situate the work within the wider agricultural-engineering challenge. The
Dalhousie University / McCain Foods project context provided the practical motivation for
investigating field-rock detection, burial-status classification and the constraints that
arise when a computer-vision model must ultimately operate close to the point of data capture.

The group acknowledges the University of Bristol Department of Engineering Mathematics and
the MSc Data Science programme for the academic environment, computing access and supporting
resources used during the project. Finally, the authors recognise the collaborative effort of
all four group members across dataset preparation, model development, conversion, Raspberry Pi
benchmarking, evidence checking, analysis and report preparation. Open sharing of experimental
outputs and careful cross-checking of results were essential to producing a coherent collective
submission.

% =============================================================================
% MAIN MATTER
% =============================================================================

\mainmatter

% =============================================================================
% CHAPTER 1
% =============================================================================

\chapter{Introduction}
\label{chap:introduction}

\section{Project Context and Motivation}

\noindent
This dissertation reports the outcomes of a four-member group project undertaken by
Tanishk Nanasaheb Shinde, Megh Kashilkar, Sanskar Sanju Gade and Jitendra Suwalka.
Unless explicitly stated otherwise, the methods, experimental results and conclusions
presented in Chapters~1--5 describe the collective project work rather than the
individual contribution of any one group member.

Precision agriculture increasingly combines sensing, computer vision and local computation
to support field operations with more timely and spatially specific information. Object
detection is particularly relevant because it can transform ordinary imagery into structured
information about the position, class and number of objects in a scene. Reviews of precision-
farming vision systems show both the rapid uptake of object detection and the continuing
importance of dataset quality, environmental variability and deployment constraints
\cite{ariza2024,badgujar2024}.

The wider project motivating this study concerns rocks in potato fields. Rocks can obstruct
field operations, contribute to equipment wear and create substantial manual removal work.
The Dalhousie University / McCain Foods project envisages detecting rocks during normal field
operations, associating detections with position information and ultimately supporting automated
removal \cite{ccrah_rock_picking,campelo2026kickoff}. The group project deliberately addresses a
narrower and verifiable part of that wider vision: detecting rocks in RGB images, distinguishing
between completely exposed and half-buried rocks, and determining whether the detector can be
run effectively on low-cost edge hardware.

\section{Problem Definition and Edge-Deployment Challenge}

The available dataset contains 448 RGB images and 706 labelled rock instances. The two target
classes are \emph{rock completely exposed} and \emph{rock half buried}. This is not only a
localisation task. Burial status matters because a partially visible rock has less visual evidence,
may share colour and texture with the surrounding soil, and is more affected by occlusion. Related
agricultural vision literature similarly identifies background complexity, changing illumination,
occlusion and limited labelled data as recurring obstacles to robust field detection
\cite{ariza2024,thurkow2025}.

A detector that performs acceptably during development is not automatically suitable for an
embedded device. The deployment representation changes the runtime, numerical precision,
memory behaviour and available optimisations even when the underlying trained model is nominally
the same. The Raspberry Pi 5 therefore creates a system-level engineering question: how much
accuracy is retained when latency is reduced, what resource profile accompanies each runtime,
and can the configuration operate reliably under repeated inference?

The project compared three lightweight candidate detectors---YOLOv8n, YOLO11n and YOLO11s---
and selected YOLOv8n using validation mAP50--95 as the predefined primary criterion. The selected
model was then represented in six edge formats: PyTorch FP32, ONNX FP32, TFLite FP32, TFLite
INT8, NCNN FP32 and NCNN FP16. The study consequently separates two decisions that can otherwise
be conflated: \emph{detector architecture selection} based on validation results, and
\emph{runtime representation selection} based on Raspberry Pi behaviour.

\section{Aim, Research Question and Hypothesis}

The principal aim is to investigate the feasibility and practical trade-offs of deploying a
two-class YOLOv8n agricultural rock detector on a Raspberry Pi 5 using alternative edge-runtime
representations.

\medskip
\noindent
The primary research question is:

\begin{quote}
\emph{How effectively can a two-class YOLOv8n rock detector be deployed on a Raspberry Pi 5,
and what accuracy, latency, memory and operational trade-offs arise across alternative
edge-runtime representations?}
\end{quote}

\noindent
The research hypothesis is that reduced-precision and edge-optimised representations will
substantially reduce Raspberry Pi inference latency relative to PyTorch FP32, but the magnitude
of improvement will vary by runtime and may involve losses in detection accuracy or operational
robustness.

\section{Objectives and Project Contribution}

The project objectives are to:
\begin{enumerate}
\item establish the two-class detector and evaluation context;
\item prepare six alternative edge-runtime representations of the selected detector;
\item create a reproducible Raspberry Pi evaluation workflow;
\item benchmark detection quality, latency, throughput and resource behaviour under a common workload;
\item investigate operational reliability and document runtime failures or negative results; and
\item select and justify a preferred deployment representation for the tested Raspberry Pi configuration.
\end{enumerate}

The principal contribution is therefore not a new detection architecture. It is an evidence-based
translation of a trained agricultural detector from model-development conditions to constrained
hardware, with the runtime alternatives compared using the same images and a common set of
accuracy, speed, counting, resource and operational measures.

\section{Scope and Boundaries}

The implemented scope includes dataset auditing, training-only augmentation, lightweight YOLO
comparison, selection of YOLOv8n, model conversion, stored-image inference on a Raspberry Pi 5,
rock counting from detections, and cross-runtime benchmarking. A separate CLIP ViT-B/32
crop-classification branch was explored during model development, but it was not part of the
verified Raspberry Pi benchmark \cite{radford2021}.

The study does \emph{not} claim implementation of live-camera operation, GPS integration,
verified rock size or weight estimation, depth or thermal sensing, robotic manipulation, or the
complete field system described in the wider project. The Raspberry Pi evaluation is therefore
an academic CPU-only proof of concept. This boundary is important because the conclusions concern
the feasibility of one detection component, not production readiness of an autonomous
agricultural machine.

\section{Dissertation Organisation}

Chapter~2 reviews the technical and agricultural background. Chapter~3 describes the executed
workflow, from dataset preparation and detector selection to Raspberry Pi conversion and
benchmarking. Chapter~4 critically evaluates model generalisation, class-specific weaknesses,
accuracy--efficiency trade-offs, resource behaviour and runtime stability. Chapter~5 answers the
research question, assesses the hypothesis and objectives, and prioritises the work required to
move from the present proof of concept toward a more realistic field system.

% =============================================================================
% CHAPTER 2
% =============================================================================

\chapter{Background and Related Work}
\label{chap:background}

\section{Precision Agriculture and Agricultural Stone Detection}

Precision agriculture applies sensing, computation and data-driven decision making at a spatial
and temporal resolution finer than conventional whole-field management. Computer vision is one
route to this objective because cameras provide a relatively inexpensive source of detailed
information about crops, weeds, obstacles and field conditions. Ariza-Sent\'is et al. review more
than one hundred recent precision-farming studies and identify object detection and tracking as
important enabling techniques, while also emphasising environmental variability and limited
labelled datasets as persistent constraints \cite{ariza2024}. A review focused specifically on
YOLO in agriculture similarly highlights the attraction of single-stage detectors for tasks that
need a practical compromise between accuracy and speed \cite{badgujar2024}.

Stone detection is a useful example of why agricultural imagery remains difficult despite strong
general-purpose vision models. Soil colour, clods, vegetation, shadows and uneven illumination
produce visually complex backgrounds, while a partly buried stone may expose only a small or
irregular region. Th\"urkow et al. demonstrate continuing research interest in agricultural stone
detection using thermal imagery, illustrating that the problem is sufficiently challenging to
motivate sensing strategies beyond conventional RGB alone \cite{thurkow2025}. The present study
retains RGB imagery because it matches the supplied project data and the intended low-cost
vision setting, but the class-specific analysis later shows why burial remains a significant source
of error.

\section{Object Detection and Lightweight YOLO Models}

Object detection combines classification with localisation. Two-stage detectors such as Faster
R-CNN first generate candidate regions and then classify/refine them, whereas one-stage detectors
predict object categories and locations more directly \cite{ren2015}. SSD demonstrated that a
single network could perform multi-scale detection without a separate proposal stage
\cite{liu2016}. YOLO pushed the single-stage idea further by framing detection as a unified
prediction problem and is now a major family of real-time detectors \cite{redmon2016}.

Modern detectors use multi-scale features because agricultural objects can occupy very different
fractions of the image. Feature Pyramid Networks formalised a widely used approach for combining
semantic information across scales \cite{lin2017}; dense detection research has also improved the
representation and optimisation of bounding boxes, including distribution-based regression in
Generalized Focal Loss \cite{li2020}. Although implementations differ between YOLO generations,
these ideas motivate the use of a compact detector that can preserve useful multi-scale features
without imposing the computation of a large backbone. The exact YOLOv8/YOLO11 implementation
used in this project is supplied by Ultralytics and should therefore be interpreted according to
the official implementation documentation rather than as a newly proposed architecture
\cite{ultralytics_docs}.

The project compared YOLOv8n, YOLO11n and YOLO11s. The ``n'' variants target low parameter and
computation budgets, while YOLO11s increases capacity. The comparison is deliberately empirical:
the architecture used for deployment was chosen by a predefined validation criterion rather than
because one generation was assumed to be universally superior.


Conceptually, the Ultralytics detectors can be separated into a backbone, a feature-fusion neck
and a detection head. The backbone converts pixels into progressively more abstract feature maps;
the neck combines information across resolutions; and the head produces class and bounding-box
predictions. Multi-scale fusion is important in this problem because a rock may occupy a large
fraction of one image but only a small region of another. YOLOv8 uses C2f blocks and an SPPF
module before multi-scale fusion, with an anchor-free decoupled detection head. The official
YOLO11 implementation changes internal modules, including C3k2/C2PSA components, while retaining
the broad lightweight one-stage design \cite{ultralytics_docs}. These details matter for model
capacity and computational demand, but they were not modified by the project; the experimental
question is which supplied lightweight model generalised best under the common training protocol.

At inference time, YOLO converts an input image into a set of candidate detections in a single
forward pass. The backbone extracts visual features, the neck propagates and combines information
at several spatial resolutions, and the decoupled head predicts bounding-box geometry and class
confidence. The resulting candidates are filtered by a confidence threshold and overlapping boxes
are suppressed during post-processing. This sequence produces the box coordinates, burial-status
class and confidence value required for annotation, class-specific counting and structured output.

The joint treatment of localisation and classification is particularly important for the present
task. A half-buried rock provides less visible shape and texture than a completely exposed rock, so
the detector must first preserve enough evidence to localise the object and then assign the correct
burial class. A missed detection, an inaccurate box and a burial-status confusion therefore represent
different failure modes even when they affect the same aggregate score. Reporting precision, recall,
F1, mAP50 and mAP50--95 gives complementary evidence about these behaviours rather than treating
``rock detected'' as a single binary outcome.

YOLO is consequently the central analytical component of the project rather than only a convenient
preprocessing tool. The same trained detector is carried from development into each edge-runtime
representation, allowing the dissertation to distinguish model-level weaknesses from changes caused
by export format, numerical precision or backend execution. This controlled continuity is necessary
for interpreting whether a deployment trade-off originates in the detector or in the runtime used to
execute it.

The ``nano'' versus ``small'' distinction is also relevant to edge deployment. Increasing model
capacity can improve representational power but does not guarantee better generalisation on a
small specialised dataset, and it increases the work required per image. Comparing YOLO11n with
YOLO11s therefore provides evidence about scale within one generation, while comparing YOLOv8n
with YOLO11n provides evidence about architecture generation at an edge-oriented scale. The
selection rule prevents these considerations from being decided by model age or parameter count
alone.

\section{Transfer Learning and Data Augmentation}

Training a detector from a pretrained representation is especially valuable when the task-specific
dataset is modest. Transfer learning exploits features learned from larger source datasets and
then adapts them to a target problem; work by Yosinski et al. shows that learned features vary in
how transferable they are but can still substantially reduce the burden of learning from scratch
\cite{yosinski2014}. Large detection datasets such as COCO also provide a common pretraining and
evaluation context for modern detectors \cite{lin2014}.

The original project dataset is small enough that overfitting and narrow coverage are material
risks. Offline augmentation was therefore applied to the training partition only. Six photometric
transformations were used: brightness, contrast, Gaussian noise, mild blur, hue/saturation and
sharpness. Applying these transformations expanded 313 original training images to 2,191 training
images. These transformations create
1,878 augmented images and 3,563 repeated annotation instances; importantly, those instances are
augmented copies of existing labelled rocks rather than 3,563 independently collected objects.
The validation and test partitions were not augmented, preserving a more meaningful estimate of
generalisation.

Agricultural scene complexity also explains why augmentation cannot by itself solve the burial
problem. A half-buried rock can be partially occluded by soil and may have ambiguous boundaries,
so the detector must infer class from less visible evidence than for a fully exposed rock. This is
a semantic and data-distribution challenge, not merely a numerical-precision problem introduced
by edge conversion.


Field imagery combines several sources of variation that are difficult to reproduce exhaustively
with augmentation. Illumination changes the apparent colour and contrast of both soil and stone;
wet and dry soil can have different reflectance; crop residues and emerging vegetation introduce
new textures; shadows can create rock-like boundaries; and camera height or viewing angle changes
the apparent scale of an object. Occlusion is especially important for burial status because the
visible fraction is itself part of the class definition. A detector may therefore correctly infer
that an object is a rock while still confusing whether it is exposed or half buried. This
motivates reporting both localisation-oriented mAP and class-specific behaviour rather than
reducing performance to a single overall score.

The class totals in the original dataset are reasonably balanced globally (362 exposed and 344
half buried), but numerical balance does not guarantee equal visual difficulty. Half-buried
instances can contain fewer discriminative pixels and more ambiguous boundaries. Consequently,
future data collection should be stratified not only by class count but by burial fraction,
lighting, soil type, scene density and rock appearance. This is a stronger strategy than simply
creating more transformed copies of the existing scenes.

\section{Detection and Counting Metrics}

For a predicted box $B_{\mathrm{pred}}$ and ground-truth box $B_{\mathrm{gt}}$, Intersection over
Union is
\begin{equation}
\mathrm{IoU}=\frac{|B_{\mathrm{pred}}\cap B_{\mathrm{gt}}|}
{|B_{\mathrm{pred}}\cup B_{\mathrm{gt}}|}.
\end{equation}
A prediction is treated as a correct detection only when its class and overlap satisfy the chosen
evaluation rules. Precision, recall and F1 are
\begin{equation}
P=\frac{TP}{TP+FP},\qquad
R=\frac{TP}{TP+FN},\qquad
F1=\frac{2PR}{P+R}.
\end{equation}
Average Precision integrates the precision--recall relationship for a class. mAP50 averages AP
across classes at IoU 0.50, while mAP50--95 averages over IoU thresholds from 0.50 to 0.95 in
steps of 0.05. The latter is stricter because it increasingly rewards well-localised boxes.

The project also evaluates counting, since a field workflow ultimately needs the number of rocks
as well as their boxes. For $N$ images, count mean absolute error is
\begin{equation}
\mathrm{MAE}=\frac{1}{N}\sum_{i=1}^{N}|\hat{y}_i-y_i|,
\end{equation}
where $y_i$ and $\hat{y}_i$ are ground-truth and predicted rock counts. Count MAE is deliberately
reported alongside mAP rather than used as a replacement: a detector may produce a similar count
while assigning an incorrect burial-status class or poorly localising the object.


Metric interpretation also depends on the operating point. A confidence threshold trades recall
against false positives, while non-maximum suppression and the IoU setting affect which overlapping
candidate boxes survive. Consequently, one confusion matrix or count error is not an intrinsic
property of a model independent of threshold choice. The benchmark fixes these settings so that
representations are compared consistently; class-specific confusion results are then used to
explain failure modes rather than to replace threshold-integrated AP measures.

\section{Edge Inference, Numerical Precision and Quantisation}

Edge computing moves processing closer to the sensor rather than depending on remote cloud
inference. In agriculture this can reduce communication latency and dependence on field
connectivity. Kalyani and Collier identify lower latency and local processing as important
motivations for cloud/fog/edge combinations in smart agriculture \cite{kalyani2021}. More recent
reviews emphasise both the opportunity for local intelligent machinery and the need for stronger
measurement of accuracy, latency, memory and resource trade-offs \cite{sapna2026}. Raspberry Pi
systems are already used across multiple precision-agriculture applications, supporting their role
as a low-cost experimentation platform while also leaving scope for more robust deployment
research \cite{joice2025}.

Model representation strongly influences edge behaviour. FP32 provides a conventional
floating-point baseline. FP16 reduces numerical storage and bandwidth, but its speed benefit
depends on whether the target runtime and processor can exploit half precision efficiently. INT8
post-training quantisation maps real values to a restricted integer range. A common affine form is
\begin{equation}
q=\operatorname{round}\left(\frac{r}{S}\right)+Z,
\end{equation}
where $r$ is the real value, $q$ the quantised integer, $S$ a scale and $Z$ a zero point. Efficient
integer-only inference can substantially reduce computation and memory traffic, but quantisation
introduces approximation error \cite{jacob2018}. Representative calibration data are therefore
important when ranges are estimated after training.


Post-training quantisation is attractive because it avoids retraining the detector, but the
calibration set becomes part of the deployment methodology. If representative activation ranges
are poorly estimated, clipping or coarse scaling can disproportionately affect difficult examples.
The project therefore drew INT8 calibration images only from original training data and balanced
the calibration object counts across the two classes. This does not guarantee lossless
quantisation, but it prevents the more serious methodological error of calibrating directly on the
final test images.

FP16 has a different trade-off. It reduces storage and numerical bandwidth without constraining
values to an eight-bit integer grid, so its accuracy is often closer to FP32. However, theoretical
savings translate into latency only when the hardware and runtime exploit half-precision
operations effectively. The later NCNN result---roughly half the model storage but only a small
latency improvement over NCNN FP32---illustrates why numerical format and realised device speed
must be measured separately.

\section{Deployment Runtimes and Inference Benchmarking}

The six representations in this study span four software ecosystems. PyTorch FP32 is the original
reference representation. ONNX is a portable graph format executed here with ONNX Runtime
\cite{onnxruntime_docs}. TFLite/LiteRT provides deployment-oriented FP32 and INT8 execution
\cite{litert_docs}. NCNN is a lightweight inference framework designed for mobile and embedded
CPU environments and supports both FP32 and FP16 representations \cite{ncnn_docs}.

Benchmarking must control the workload if runtime effects are to be interpreted fairly. MLPerf
Inference formalises the broader principle that inference comparisons require clear scenarios,
metrics and reproducible conditions \cite{reddi2020}. This dissertation does not implement the
MLPerf suite, but adopts the same discipline: the completed representations process the same
90 images at the same $640\times640$ input resolution and batch size one, with the same detection
threshold settings. Latency, throughput, CPU utilisation, process memory, temperature, detection
quality and count error are considered jointly.

Latency and FPS alone are insufficient for a constrained device. A backend can be fast while
using high CPU, accumulating heat or failing under repeated inference. Likewise, a highly accurate
representation may be too slow for a practical on-the-go workflow. The central analytical idea of
this dissertation is therefore an accuracy--efficiency--robustness trade-off rather than a single
leaderboard metric.


A controlled benchmark also needs a clear timing boundary. End-to-end camera latency, image
preprocessing, model execution and post-processing can all be measured, but combining different
boundaries across runtimes would make comparison misleading. The present experiment standardises
input geometry in advance and concentrates on the implemented inference path, while separately
acknowledging that a future live-camera system will introduce additional acquisition and pipeline
costs. Warm-up inference is used because the first call can include graph initialisation, memory
allocation or cache effects that are not representative of steady execution.

Mean latency describes the average timed cost, while median FPS provides an intuitive measure of
typical throughput. Resource measures add complementary information: process RSS captures memory
resident for the running process; CPU utilisation indicates how aggressively the backend uses the
processor; and temperature provides a proxy for accumulated thermal load. None is sufficient on
its own, which is why the later selection considers the joint evidence.

\section{Research Gap}

The literature establishes three connected points. First, agricultural object detection is useful
but sensitive to field variability \cite{ariza2024,badgujar2024}. Second, local edge processing is
attractive where latency and connectivity matter, yet deployment studies need multidimensional
resource evaluation \cite{kalyani2021,sapna2026}. Third, Raspberry Pi platforms are sufficiently
capable to support precision-agriculture applications but remain constrained systems
\cite{joice2025,raspberrypi5_docs}.

The gap addressed here is consequently not whether YOLO can detect agricultural objects in
principle. It is the lack of evidence that a detector selected in a development environment will
retain an appropriate combination of accuracy, latency, memory use and operational behaviour
when moved into alternative runtime representations on low-cost hardware. Holding the detector,
images and target device constant while changing the representation allows the project to isolate
this practical deployment decision.


This position also distinguishes the dissertation from studies that report only compressed model
size or desktop inference speed. A format that is theoretically portable may still fail because of
software compatibility, memory pressure or device-level instability, while a smaller numeric
representation may provide little acceleration if the target CPU cannot exploit it. Conversely, a
runtime can provide large speed gains without changing the learned weights materially. Measuring
all formats on the Raspberry Pi is therefore necessary to convert general expectations about
quantisation and lightweight inference into device-specific evidence.

The review consequently motivates two linked evaluations. The first asks whether the selected
detector generalises from validation to unseen project images and whether both burial classes are
recognised comparably. The second asks how much of that detection behaviour survives deployment
and what speed/resource consequences accompany each representation. Chapters~3 and~4 preserve
this separation so that an architecture weakness is not incorrectly blamed on a runtime, and a
runtime failure is not incorrectly interpreted as evidence that the detector architecture itself
is invalid.


A final background issue is the distinction between model size, process memory and computational
cost. A smaller model file reduces storage and transfer requirements, but resident memory also
contains runtime libraries and working buffers; similarly, fewer bytes per parameter do not imply
a proportional reduction in latency. FLOP-style complexity estimates are useful for architecture
comparison, yet they do not capture operator implementation, thread scheduling or memory access on
a specific CPU. For that reason, the dissertation treats file/model representation as one factor
and relies on measured RSS, CPU utilisation and latency for the deployment decision.


Reproducibility also requires negative outcomes to remain visible. If a runtime exports successfully
but resets the device during a longer workload, excluding it without explanation would bias the
comparison toward successful formats. Equally, filling the missing row with estimates from a few
images would make the table appear complete at the expense of validity. The methodological position
adopted here is therefore conservative: only completed common-workload runs receive aggregate metrics,
and incomplete or failed runs are reported qualitatively with their evidential limits. This principle is
particularly important for embedded systems, where software compatibility and operational stability are
part of deployment performance rather than peripheral implementation details.

% =============================================================================
% CHAPTER 3
% =============================================================================

\chapter{Design and Implementation}
\label{chap:execution}

\section{Implementation Overview}

The implementation consisted of two connected stages. The first established a two-class detector
for identifying rocks in agricultural RGB images and distinguishing completely exposed from
half-buried rocks. The second transferred the selected detector to a Raspberry Pi 5 and evaluated
alternative representations under a controlled stored-image workload.

\section{Dataset Preparation and Detector Selection}

The original dataset contained 448 images and 706 annotated rocks. Before augmentation it was
partitioned into training, validation and test sets as shown in Table~\ref{tab:dataset}. The test
set contained 90 images and 135 labelled rocks and was not used to select the detector architecture.

\begin{table}[htbp]
\centering
\caption{Dataset split and class composition before augmentation.}
\label{tab:dataset}
\begin{tabular}{|l|r|r|r|}
\hline
Split & Images & Completely exposed & Half buried \\
\hline
Training & 313 & 259 & 250 \\
\hline
Validation & 45 & 34 & 28 \\
\hline
Test & 90 & 69 & 66 \\
\hline
Total & 448 & 362 & 344 \\
\hline
\end{tabular}
\end{table}

\reportchart[0.76\textwidth]{figures/figure_01.png}{Distribution of the 448 original images across the training, validation and test partitions. The fixed pre-augmentation split separates detector fitting from architecture selection and the later 90-image evaluation workload.}{fig:dataset-split-chart}

\reportchart[0.76\textwidth]{figures/figure_02.png}{Class-level annotation counts for completely exposed and half-buried rocks in each dataset split. The closely balanced overall totals allow later performance differences to be interpreted primarily as differences in visual difficulty rather than severe label imbalance.}{fig:class-distribution-chart}

\FloatBarrier

Only the training set was expanded offline. Each of the 313 training images received six
photometric variants, producing 1,878 augmented images in addition to the originals. The resulting
2,191-image training set contained 3,563 annotation instances across augmented copies. The
augmentation strategy targeted plausible appearance variation while avoiding geometric changes
that could undermine the existing boxes.

\reportchart[0.76\textwidth]{figures/figure_03.png}{Expansion of the training partition from 313 original images to 2,191 images after six photometric variants were generated per source image. Only training data were augmented; validation and test images remained unchanged.}{fig:training-augmentation-chart}

\reportchart[0.76\textwidth]{figures/figure_04.png}{Dataset-scale summary showing 448 original images, 706 annotated rock objects and two burial-status classes. These figures define the limited-data setting that motivates transfer learning, controlled augmentation and careful held-out evaluation.}{fig:dataset-summary-chart}

\FloatBarrier

YOLOv8n, YOLO11n and YOLO11s were trained under a common main configuration: a maximum of
20 epochs, image size 896, batch size 32, seed 42, automatic mixed precision, cosine learning-rate
scheduling and Ultralytics automatic optimiser selection, which selected AdamW. Mosaic was
disabled; horizontal flip probability was 0.5, translation 0.04 and scale 0.1. Final comparative
validation inference used 1024-pixel images.


The common training settings were intended to make the architecture comparison interpretable. The
candidate models were not given independently tuned hyperparameters that could favour one model.
Training curves and validation outputs were retained for each candidate, and the final selection
was made only after the common validation inference. This design is less exhaustive than a full
hyperparameter search, but it provides a controlled comparison appropriate to the project scope.

Table~\ref{tab:modelselection} reports the validation comparison. YOLO11s achieved the best F1,
recall and mAP50, but YOLOv8n achieved the highest mAP50--95. Because mAP50--95 was defined in
advance as the primary architecture-selection criterion, YOLOv8n was selected. This avoids a
post-hoc claim that the chosen model was ``best'' on every measure.

\begin{table}[htbp]
\centering
\caption{Validation comparison used for detector selection.}
\label{tab:modelselection}
\begin{tabular}{|l|c|c|c|c|c|}
\hline
Model & Precision & Recall & F1 & mAP50 & mAP50--95 \\
\hline
YOLOv8n  & 0.681 & 0.759 & 0.718 & 0.759 & \textbf{0.4288} \\
\hline
YOLO11n  & 0.743 & 0.719 & 0.731 & 0.729 & 0.4002 \\
\hline
YOLO11s  & 0.764 & \textbf{0.801} & \textbf{0.782} & \textbf{0.783} & 0.4107 \\
\hline
\end{tabular}
\end{table}

\reportchart[0.76\textwidth]{figures/figure_05.png}{Validation and test precision, recall and F1 for YOLOv8n, YOLO11n and YOLO11s. YOLO11s leads several operating-point metrics, showing why the detector choice cannot be justified from a single score in isolation.}{fig:precision-recall-f1-chart}

\reportchart[0.76\textwidth]{figures/figure_06.png}{mAP50 and mAP50--95 for the three candidate YOLO detectors on validation and test data. The stricter mAP50--95 measure evaluates localisation across multiple IoU thresholds and was the predefined architecture-selection criterion.}{fig:map-comparison-chart}

\begin{figure}[!htbp]
\centering
\begin{minipage}[t]{0.315\textwidth}
  \centering
  \includegraphics[width=\linewidth,keepaspectratio]{figures/figure_07_model_size.png}
  \vspace{0.25em}
  \raggedright\footnotesize
  \textbf{(a) Stored model size.} YOLO11n has the smallest stored footprint at
  5.21~MB, YOLOv8n remains compact at 6.00~MB, and YOLO11s is substantially
  larger at 18.35~MB.
\end{minipage}\hfill
\begin{minipage}[t]{0.315\textwidth}
  \centering
  \includegraphics[width=\linewidth,keepaspectratio]{figures/figure_07_inference_speed.png}
  \vspace{0.25em}
  \raggedright\footnotesize
  \textbf{(b) Validation inference speed.} Development-side throughput is
  nearly identical across the candidates, with YOLOv8n marginally highest at
  8.53~FPS.
\end{minipage}\hfill
\begin{minipage}[t]{0.315\textwidth}
  \centering
  \includegraphics[width=\linewidth,keepaspectratio]{figures/figure_07_validation_map.png}
  \vspace{0.25em}
  \raggedright\footnotesize
  \textbf{(c) Validation localisation quality.} YOLOv8n records the strongest
  validation mAP50--95 at 0.4288, the predefined detector-selection criterion,
  ahead of YOLO11s and YOLO11n.
\end{minipage}
\caption{YOLO candidate efficiency and validation quality.}
\label{fig:model-efficiency-chart}
\end{figure}

\FloatBarrier

Evaluation of the selected YOLOv8n on the untouched architecture-selection test set produced
mAP50--95 of approximately 0.3384, compared with 0.4288 on validation. The absolute reduction
of 0.0904 (about 21\% relative) indicated a meaningful generalisation gap and motivated caution
when interpreting all later deployment results.


The ordering of these steps is important. Architecture choice was fixed from validation before the
test value was examined, so the 90-image partition retained its intended role for estimating the
selected architecture at that stage. Once the same 90 images were later reused to compare edge
representations and choose a runtime, however, they ceased to be a completely untouched final
deployment test. The dissertation therefore distinguishes the validity of the original
architecture-selection use from the later limitation introduced by runtime selection.

The model-development records also retained class-level outputs, confusion information and saved
weights rather than only a single best score. This made it possible to identify the half-buried
weakness before deployment. Without that evidence, a lower class performance observed after INT8
conversion could easily have been misattributed to quantisation rather than recognised as a
pre-existing detector limitation.

A separate CLIP ViT-B/32 branch investigated whether pretrained vision--language embeddings
could distinguish burial status from rock crops \cite{radford2021}. This work remained a
supporting model-development branch. No verified Raspberry Pi CLIP benchmark was completed,
and the deployment experiments below concern YOLOv8n only.


The research branch used CLIP ViT-B/32 in two modes: zero-shot classification using textual class
descriptions, and a lightweight linear probe trained on frozen image embeddings. A confidence-
based rule was explored for disagreements between the YOLO class and CLIP classification, with a
configured CLIP override threshold of 0.80. Conceptually this created a pipeline of
\emph{field image $\rightarrow$ YOLO detection $\rightarrow$ rock crop $\rightarrow$ CLIP
verification}. The branch was intentionally excluded from the final Pi comparison because running
CLIP on every detected crop would add a second model, additional memory demand and extra latency,
and no complete device benchmark existed to quantify that cost.

\section{YOLO Detection and Output Processing}

The selected YOLOv8n detector is the primary executable model in the final system. For each input
image it produces a collection of bounding boxes, confidence values and class identifiers in one
forward pass. During development, the common training and validation configuration was used to
compare candidate architectures. During edge evaluation, the selected weights received the same
fixed $640\times640$ letterboxed pixels in every runtime so that the comparison did not confound
detector behaviour with different image geometry.

Post-processing used the common confidence threshold of 0.25 and IoU setting of 0.70. Candidate
boxes were converted into the common coordinate system, associated with either the completely
exposed or half-buried class, and retained in structured per-image outputs. These predictions drove
the annotated images, class-specific counts and the precision, recall, F1 and mAP calculations.
Treating the output as structured detection evidence, rather than only a rendered image, made it
possible to trace aggregate results back to individual images and class-specific errors.

The two-class detection formulation also avoids requiring a second classifier for the deployed
benchmark: localisation and burial status are produced by the same compact network. This reduces
pipeline complexity and keeps runtime comparison focused on one trained model. The trade-off is
that weak visual evidence for a half-buried rock can affect both localisation and class assignment,
which is why the report gives the detector-selection evidence and class-specific behaviour greater
weight than an FPS-only deployment comparison.

Holding the YOLOv8n weights, input pixels, thresholds and output interpretation constant across
PyTorch, TFLite and NCNN is the principal experimental control. Observed changes can therefore be
analysed as representation and runtime effects while retaining the same underlying detection task.

\section{Edge-Deployment Design and Test Workload}

The selected YOLOv8n model was prepared in six representations: PyTorch FP32, ONNX FP32,
TFLite FP32, TFLite INT8, NCNN FP32 and NCNN FP16. These formats were chosen to span the
original framework, an exchange/runtime path, an integer-quantised deployment path and a
lightweight embedded inference framework.


All six variants originated from the same selected YOLOv8n weights. They were not separately
retrained after export. This is a key control in the experiment: when a converted representation
changes accuracy or latency, the difference is primarily associated with conversion, numerical
precision and backend execution rather than a different optimisation history. Export success was
also separated from deployment success. A file existing on disk or loading once was not sufficient
to classify a representation as benchmark-complete.

For the edge comparison, the same 90 test images were converted into fixed $640\times640$
letterboxed images. If the original dimensions are $W\times H$, the aspect-ratio-preserving
scale is
\begin{equation}
s=\min\left(\frac{640}{W},\frac{640}{H}\right),
\end{equation}
with symmetric padding added using fill value 114. This produces a common input shape while
avoiding aspect-ratio distortion. The 90 images contain 135 labelled rocks: 69 completely exposed
and 66 half buried.


Bounding-box centres and dimensions were transformed using the same scale and padding offsets as
the image. A manifest retained original dimensions, deployment dimensions, scaling and padding
information so that the fixed images could be traced back to their sources. Pre-generating the
letterboxed workload reduced framework-specific preprocessing as a confounding variable: each
backend received the same pixels at the same geometry instead of relying on slightly different
runtime resize implementations.

A ten-image smoke-test subset was copied from these same 90 images to verify model loading,
input/output handling and basic post-processing before longer runs. It is therefore a functional
subset, not an independent evaluation dataset.


The smoke-test procedure checked model loading, image acceptance, inference, class identifier
mapping, bounding-box generation, class-specific counting, annotated-image output and structured
JSON output. In the final functional check, the optimised formats returned the expected single
completely exposed rock on the representative image. Passing this test established interface
compatibility only; it did not contribute additional accuracy samples beyond the 90-image set.

The INT8 conversion used 200 \emph{original training images} as representative calibration data,
not validation or test images. The selected calibration set contained 304 labelled objects,
balanced at 152 completely exposed and 152 half buried. This separation prevents direct leakage
of the final test images into calibration while ensuring that both classes influence activation
range estimation.


The calibration set contained 304 labelled objects in total. Equal object counts by class do not
mean the calibration data are statistically identical to the test distribution, but the design
ensures that one burial class does not dominate simply because it contributes more representative
objects. After conversion, the calibration images themselves were not required for device
inference; provenance was retained through the calibration manifest rather than by treating the
calibration subset as a second evaluation set.

\section{Raspberry Pi Environment and Tooling}

The target device was a Raspberry Pi 5 Model B with a 64-bit ARM environment and approximately
8~GB RAM \cite{raspberrypi5_docs}. It ran Debian GNU/Linux 13 and Python 3.13.5. Key package
versions were Ultralytics 8.4.115, ONNX Runtime 1.28.0, OpenCV 5.0.0, NumPy 2.5.1,
Pandas 3.0.5 and psutil 7.2.2. LiteRT inference used the
\texttt{ai\_edge\_litert.interpreter} interface. NCNN was used for the FP32 and FP16 variants.

Power was supplied by the official Raspberry Pi 27~W USB-C supply rated at 5.1~V/5~A. No Active
Cooler was available in the final configuration; Raspberry Pi documentation treats cooling as a
material consideration for sustained high-load operation \cite{raspberrypi_cooling}. No live
camera was used. A MacBook served only for SSH/SCP control and file transfer; measured inference
ran on the Pi CPU.


\begin{table}[htbp]
\centering
\caption{Key Raspberry Pi benchmark environment.}
\label{tab:environment}
\begin{tabular}{|p{5.0cm}|p{8.4cm}|}
\hline
Component & Recorded configuration \\
\hline
Device & Raspberry Pi 5, aarch64, approximately 8~GB RAM \\
\hline
Operating system & Debian GNU/Linux 13 \\
\hline
Python / Ultralytics & Python 3.13.5 / Ultralytics 8.4.115 \\
\hline
ONNX Runtime & 1.28.0 \\
\hline
OpenCV / NumPy / Pandas & 5.0.0 / 2.5.1 / 3.0.5 \\
\hline
Resource monitoring & psutil 7.2.2 \\
\hline
LiteRT interface & \texttt{ai\_edge\_litert.interpreter} \\
\hline
Power & Official Raspberry Pi 27~W, 5.1~V/5~A \\
\hline
Cooling & No Active Cooler \\
\hline
Input source & Stored images; no camera \\
\hline
\end{tabular}
\end{table}

Recording exact versions is important because inference behaviour is not determined by the model
file alone. Runtime releases, Python bindings, numerical libraries and operating-system changes can
alter performance or compatibility. The environment table therefore forms part of the
reproducibility evidence rather than merely documenting convenience software.

Deployment scripts performed model-specific loading, image reading, inference, post-processing,
resource sampling and CSV/JSON output. Output manifests and saved predictions were used to verify
that each completed run contained the intended 90 images rather than silently omitting failures.


The tooling was separated by purpose. Environment verification recorded software versions before
benchmarking; single-image inference provided a low-cost functional test; the benchmark script
produced per-image timing/resource records and aggregate outputs; and a dedicated stability script
performed repeated continuous inference. SSH and SCP scripts supported transfer between the
MacBook and Pi without moving the actual inference workload off-device. This separation reduced
manual intervention and made it easier to distinguish an environment problem, a model-interface
problem and a sustained-runtime problem.

For timing, model loading was kept outside the principal per-image inference measurement where
possible, and warm-up calls were used before the timed pass. CPU and process RSS were sampled as
system-level indicators rather than interpreted as exact energy measurements. Temperature was
recorded from the Raspberry Pi monitoring interface. These measurements therefore characterise
the implemented software stack on the Pi, not an abstract hardware-only cost of the neural
network.

\section{Benchmark and Evaluation Protocol}

The final executed benchmark used batch size one, $640\times640$ inputs, confidence threshold
0.25 and IoU setting 0.70. Each representation received three warm-up inferences followed by one
timed pass across the 90-image workload. Table~\ref{tab:protocol} distinguishes this from an earlier
planning document that proposed ten warm-ups, three timed passes and a 30-minute stability test.
Those planned settings were not executed and are not used as if they were observed evidence.

\begin{table}[htbp]
\centering
\caption{Planned protocol versus the protocol actually executed.}
\label{tab:protocol}
\begin{tabular}{|p{4.2cm}|p{4.8cm}|p{4.8cm}|}
\hline
Element & Earlier plan & Executed benchmark \\
\hline
Input size & $640\times640$ & $640\times640$ \\
\hline
Batch size & 1 & 1 \\
\hline
Warm-up & 10 inferences & 3 inferences \\
\hline
Timed evaluation & 3 passes & 1 pass \\
\hline
Images per completed format & 90 per pass & 90 \\
\hline
Stability target & 30 minutes & Separate short tests only \\
\hline
Completed formats & Planned six & Five complete; ONNX aggregate N/A \\
\hline
\end{tabular}
\end{table}

For each completed representation, the benchmark recorded inference latency, derived throughput,
CPU utilisation, process resident memory, temperature and predictions. Detection outputs were then
evaluated against the common labels to obtain precision, recall, F1, mAP50, mAP50--95 and count
MAE. Median FPS is reported consistently in the final comparison because it is less sensitive to
isolated timing spikes than an arithmetic mean of reciprocal latencies.


Latency was recorded per image after warm-up, allowing both aggregate means and the distribution of
individual timings to be inspected. Throughput was derived from the observed inference rate rather
than from a theoretical operation count. Process memory was tracked as RSS, which captures the
resident memory of the running process but not every system allocation attributable indirectly to
the workload. CPU percentage is likewise a software/runtime measure: a backend can achieve lower
latency by using more of the available cores, so high CPU utilisation is interpreted together with
speed rather than automatically penalised.

The prediction record retained class labels, confidences and bounding boxes for downstream
accuracy validation. Rock count was obtained from the retained detections per image, making count
MAE directly traceable to the same predictions used for mAP. Temperature and throttle indicators
were captured during benchmark/stability runs so that a high FPS result could be checked for signs
of unsustainable operating conditions.


The accuracy evaluation used the predictions saved by each completed representation rather than
mixing benchmark timings from one run with detections from a different workload. This allowed the
same labels to support both detection metrics and per-image rock-count error. It also made failure
visible: a representation that did not produce the complete prediction set could not be assigned
an aggregate row by extrapolating from a small subset.

\section{Reliability, Diagnostics and Reproducibility}

Five representations completed the 90-image benchmark: PyTorch FP32, TFLite FP32, TFLite INT8,
NCNN FP32 and NCNN FP16. ONNX FP32 exported correctly and processed isolated images, but repeated
sustained attempts ended in abrupt Raspberry Pi reboots before a valid aggregate run was produced.
The root cause was not established, so no ONNX aggregate latency, accuracy or resource values are
estimated. The defensible conclusion is only that ONNX FP32 was functionally viable for isolated
inference but did not demonstrate stable sustained benchmarking in the tested configuration.


Diagnostic logs were reviewed for evidence that could distinguish an application exception from a
system reset. The available records established the temporal association between sustained ONNX
attempts and abrupt reboot, but did not uniquely identify thermal throttling, undervoltage,
out-of-memory termination or an ONNX Runtime defect. This is why the final benchmark table marks
ONNX as N/A instead of reporting values from a partial run. Preserving an incomplete row is less
visually satisfying than a six-way comparison, but avoids false precision.

Operational testing also exposed the difference between completing a short fixed workload and
remaining stable continuously. NCNN FP32 completed the 90-image benchmark but a later continuous
run hard-reset the Pi and did not produce a valid completed continuous CSV. By contrast, TFLite
INT8 completed a genuine continuous experiment lasting approximately two minutes, producing
about 2,533 inference rows at roughly 22.7 FPS. Temperature increased from the mid-40s to about
76--77~$^\circ$C, with a raw recorded maximum near 77.1~$^\circ$C, and no throttling or reboot was
recorded during that run.

This is deliberately described as short-duration evidence. No valid 10-minute or 30-minute final
stability run exists. A file whose name suggested a 30-minute test was found to duplicate the
shorter dataset and is not treated as independent evidence. This distinction between planned,
completed and invalid evidence is part of the reproducibility approach rather than a limitation
to be hidden.


Reproducibility was supported by retaining the deployment folder structure, environment records,
input manifests, prediction JSON files, annotated outputs and benchmark CSVs. The evidence was
also checked against duplicate backups so that copied result folders were not counted as separate
experiments. Where a long-run file conflicted with its contents, the contents and row counts took
precedence over the filename. This evidence-first rule is why the dissertation reports one genuine
approximately two-minute INT8 stability run and no completed 10- or 30-minute final run.


The same principle was applied to backup folders. Copies of a benchmark were treated as backups,
not additional experimental repetitions. Annotated-image and prediction-JSON counts were used to
confirm the expected 90-image completion for formats such as NCNN FP16. This prevents the number
of files in an archive from being mistaken for the number of independent trials.

\section{Chapter Summary}

The implementation created a controlled path from an empirically selected YOLOv8n detector to
five fully benchmarked Raspberry Pi representations plus an ONNX format that remained incomplete
under sustained testing. The next chapter evaluates what those measurements mean, with particular
attention to the pre-existing half-buried weakness and the trade-off between PyTorch accuracy,
INT8 speed and NCNN FP16 accuracy retention.

% =============================================================================
% CHAPTER 4
% =============================================================================

\chapter{Critical Evaluation}
\label{chap:evaluation}

\section{Evaluation Overview}

The central evaluation result is that deployment quality cannot be represented by a single metric.
The representation with the strongest measured detection quality was the slowest, while the fastest
representation introduced a measurable accuracy penalty. Resource utilisation and operational
behaviour also differed between backends. Table~\ref{tab:benchmark} therefore presents the core
cross-format evidence in one place.

\begin{table}[htbp]
\centering
\caption{Canonical Raspberry Pi benchmark results. FPS denotes median throughput.}
\label{tab:benchmark}
\resizebox{\textwidth}{!}{%
\begin{tabular}{|l|r|r|r|r|r|r|r|r|r|}
\hline
Representation & Latency ms & FPS & CPU \% & Peak RAM MB & P & R & F1 & mAP50 & mAP50--95 \\
\hline
PyTorch FP32 & 397.48 & 2.52 & 44.47 & 839.47 & 0.716 & 0.550 & 0.622 & 0.548 & 0.275 \\
\hline
TFLite FP32  & 143.70 & 6.95 & 69.95 & 868.27 & 0.669 & 0.542 & 0.599 & 0.535 & 0.270 \\
\hline
TFLite INT8  & 44.13  & 22.62 & 62.83 & 824.22 & 0.661 & 0.539 & 0.594 & 0.503 & 0.234 \\
\hline
NCNN FP32    & 75.09  & 13.26 & 83.32 & 835.25 & 0.667 & 0.542 & 0.598 & 0.535 & 0.271 \\
\hline
NCNN FP16    & 72.31  & 13.82 & 87.08 & 835.45 & 0.672 & 0.542 & 0.600 & 0.537 & 0.270 \\
\hline
ONNX FP32    & N/A & N/A & N/A & N/A & N/A & N/A & N/A & N/A & N/A \\
\hline
\end{tabular}}
\end{table}

\begin{figure}[!htbp]
\centering
\begin{minipage}{0.96\textwidth}
  \centering
  \includegraphics[width=\linewidth,keepaspectratio]{figures/figure_table_4_1_runtime.png}
  \vspace{0.15em}
  \raggedright\footnotesize\textbf{(a) Runtime performance:} mean inference
  latency and median throughput for the five completed representations.
\end{minipage}

\vspace{0.35em}
\begin{minipage}{0.96\textwidth}
  \centering
  \includegraphics[width=\linewidth,keepaspectratio]{figures/figure_table_4_1_resources.png}
  \vspace{0.15em}
  \raggedright\footnotesize\textbf{(b) Resource demand:} mean CPU usage and
  peak process memory measured during the same 90-image benchmark.
\end{minipage}

\vspace{0.35em}
\begin{minipage}{0.96\textwidth}
  \centering
  \includegraphics[width=\linewidth,keepaspectratio]{figures/figure_table_4_1_accuracy.png}
  \vspace{0.15em}
  \raggedright\footnotesize\textbf{(c) Detection quality:} precision, recall,
  F1, mAP50 and mAP50--95 retained after deployment conversion.
\end{minipage}

\vspace{0.45em}
\begin{minipage}{0.96\textwidth}
\footnotesize
\textbf{Part 1 -- Efficiency and resources.} Images (a) and (b) compare the
runtime trade-off across latency, throughput, CPU usage and peak process memory;
TFLite INT8 provides the strongest speed and lowest memory result, while the
NCNN formats require the greatest CPU share.

\smallskip
\textbf{Part 2 -- Detection quality.} Image (c) separates operating-point
precision, recall and F1 from the stricter mAP measures. PyTorch FP32 retains
the strongest aggregate accuracy, whereas TFLite INT8 exchanges some
localisation quality for substantially faster inference. ONNX FP32 is omitted
because no valid aggregate benchmark was completed.
\end{minipage}
\caption{Three-image graphical summary of Table~\ref{tab:benchmark}.}
\label{fig:benchmark-summary-chart}
\end{figure}

\FloatBarrier

Count MAE was 0.600 for PyTorch FP32, TFLite FP32 and NCNN FP32, 0.611 for NCNN FP16 and
0.633 for TFLite INT8. The close values indicate that quantisation did not create a catastrophic
counting failure, but the modestly higher INT8 error is consistent with its lower detection metrics.

\section{Detector Selection and Class-Specific Behaviour}

The detector-selection result should be interpreted according to the stated criterion. YOLO11s
had the highest validation mAP50 (0.783), recall (0.801) and F1 (0.782), whereas YOLOv8n had the
highest validation mAP50--95 (0.4288). Selecting YOLOv8n was therefore internally valid because
mAP50--95 had been defined as the primary criterion; it would be incorrect to claim that YOLOv8n
was best across every metric.

The untouched test result also shows why validation leadership should not be mistaken for final
generalisation. YOLOv8n mAP50--95 decreased from 0.4288 on validation to approximately 0.3384
on the untouched test set. A 0.0904 absolute drop, around 21\% relative, suggests that the modest
dataset and heterogeneous field appearance leave material uncertainty beyond the development
partition. This is consistent with literature that identifies generalisation under changing
agricultural conditions as a persistent challenge \cite{ariza2024,badgujar2024}.


Several mechanisms could contribute to this gap: the small validation partition may provide a
noisy estimate, the held-out images may contain harder scenes, or the selected training recipe may
not cover the full variation in soil, illumination and occlusion. The experiment does not isolate
which explanation dominates. The appropriate conclusion is therefore that the validation result
was optimistic relative to the held-out test, not that a specific data defect or optimisation
choice was proven to cause the reduction.

The largest semantic weakness was class-specific. At the evaluated source-detector operating
point, 64 of 69 completely exposed rocks were correctly represented in the relevant confusion
result (about 92.8\%), compared with only 23 of 66 half-buried rocks (about 34.8\%). These are
operating-point correct rates from the confusion analysis, not AP values. Their importance is that
the burial problem existed before Raspberry Pi conversion. INT8 and FP16 can alter confidence or
localisation, but they cannot plausibly be treated as the sole cause of a weakness already visible
in the source detector.

The result has a direct design implication: future effort spent only on runtime optimisation may
improve FPS while leaving the main application-level recognition weakness substantially unchanged.
More representative half-buried examples, stronger hard-negative coverage and perhaps sensing
modalities less sensitive to soil/stone appearance are likely to produce greater field value than
further micro-optimisation of an already fast INT8 runtime.


This asymmetry also affects how aggregate mAP should be interpreted operationally. An average over
two classes can hide the fact that one class is close to the intended behaviour while the other is
frequently missed. In a rock-removal workflow, missing a half-buried stone may be more costly than
a modest localisation error on a fully exposed stone. Future evaluation should therefore report
per-class AP/recall and field-relevant miss rates alongside aggregate metrics, and should consider
whether class-specific confidence thresholds are justified.

\section{Accuracy--Efficiency Trade-offs Across Edge Formats}

Relative to the 397.48~ms PyTorch baseline, TFLite FP32 reduced mean latency by approximately
2.77$\times$, NCNN FP32 by 5.29$\times$, NCNN FP16 by 5.50$\times$, and TFLite INT8 by
9.01$\times$.

\reportchart[0.76\textwidth]{figures/figure_deployment_speed_accuracy_scatter.png}{Median throughput plotted against mAP50--95 for each deployment format that completed the Raspberry Pi benchmark. Points farther right are faster and points higher on the chart retain stronger localisation quality; ONNX FP32 is omitted because no valid aggregate benchmark was completed.}{fig:deployment-speed-accuracy-scatter}

\reportchart[0.82\textwidth]{figures/figure_09.png}{Median throughput and corresponding mean inference latency for each completed deployment format on Raspberry Pi 5. TFLite INT8 provides the highest speed at 22.62 FPS and the lowest latency at 44.13~ms.}{fig:fps-latency-chart}

\reportchart[0.82\textwidth]{figures/figure_10.png}{F1, mAP50 and mAP50--95 for the five deployment formats that completed the 90-image benchmark. PyTorch retains the strongest aggregate accuracy, while INT8 trades measurable localisation quality for substantially faster inference.}{fig:edge-accuracy-chart}

PyTorch FP32 achieved the highest edge-benchmark mAP50 at 0.548 and mAP50--95 at 0.275, but
2.52 median FPS is restrictive for an on-the-go vision component. TFLite FP32 preserved most of
that accuracy (mAP50 0.535) while raising throughput to 6.95 FPS. NCNN FP32 produced almost the
same accuracy (mAP50 0.535; mAP50--95 0.271) at 13.26 FPS, demonstrating that runtime choice
alone can yield a large speed improvement without material loss in the measured accuracy.

NCNN FP16 produced mAP50 0.537 and mAP50--95 0.270 at 13.82 FPS. Its model storage was roughly
49\% smaller than NCNN FP32, yet latency improved by only about 3.8\%. This is an important
counterexample to the assumption that halving numerical width necessarily produces a comparable
CPU speedup. On this Raspberry Pi/runtime combination, storage reduction was much larger than
execution-time reduction.


The closeness of the three non-quantised deployment results is also notable. TFLite FP32 and NCNN
FP32 both achieved mAP50 0.535, while NCNN FP16 slightly increased this to 0.537; all three were
within 0.013 mAP50 of the PyTorch baseline. Their mAP50--95 values (0.270, 0.271 and 0.270) were
similarly close to PyTorch at 0.275. Within the resolution of this experiment, exporting to these
representations preserved most of the detector's aggregate quality. That strengthens the argument
that runtime optimisation can produce meaningful speed improvements before accepting the larger
quantisation penalty of INT8.

Precision and recall changed only modestly across these three formats. PyTorch precision was
0.716 and recall 0.550; the converted FP32/FP16 formats clustered around 0.667--0.672 precision and
0.542 recall. TFLite INT8 moved to 0.661 precision and 0.539 recall. The direction is consistent
with a small degradation rather than a wholesale change in detector behaviour, but the lower
mAP50--95 indicates that stricter localisation quality is more sensitive than the headline
operating-point recall alone suggests.

TFLite INT8 delivered the largest performance change: 44.13~ms mean latency and 22.62 median
FPS, approximately nine times the PyTorch speed. The cost was a decline in mAP50 from 0.548 to
0.503 (0.045 absolute, about 8.2\%) and in mAP50--95 from 0.275 to 0.234 (0.041 absolute, about
14.9\%). This behaviour is consistent with the general quantisation trade-off described in the
literature: integer arithmetic can make inference substantially cheaper, while approximated
weights/activations can reduce task accuracy \cite{jacob2018}.

The practical question is therefore whether the lost accuracy is acceptable for the intended
stage of the system. For a proof-of-concept CPU-only edge detector, moving from 2.52 to 22.62 FPS
changes the operational character of the application. The corresponding mAP loss is real and must
not be minimised, but it is sufficiently bounded that INT8 remains a credible choice when latency
is a first-order requirement.


The FP32 runtime comparison is particularly informative because it separates numerical precision
from backend implementation. PyTorch FP32, TFLite FP32 and NCNN FP32 all use nominally full
precision, yet mean latency changes from 397.48 to 143.70 to 75.09~ms. A large part of the edge
speed-up therefore comes from the inference engine and graph/runtime optimisation before reduced
precision is introduced. INT8 adds a further major gain on top of this runtime effect. This staged
interpretation is more defensible than attributing the entire 9.01$\times$ improvement to
quantisation alone.


The result also shows why throughput should be compared using the same batch-one scenario. Large
batches can increase hardware utilisation on servers but do not represent a camera-like stream in
which each frame needs a prompt result. At 22.62 median FPS, the INT8 representation has a nominal
per-frame budget compatible with substantially more responsive processing than the 2.52 FPS
PyTorch baseline. Whether that is sufficient for a moving planter depends on camera field of view,
vehicle speed and required overlap between observations---quantities not measured here---so the
study stops short of translating FPS into a guaranteed maximum tractor speed.

Accuracy differences should likewise be treated as estimates on 90 images rather than exact
population values. The reported decimal metrics are useful for comparing completed runs on a
common set, but the experiment was not large enough to establish that a difference of a few
thousandths between TFLite FP32 and NCNN FP16 would reproduce across independent farms. The broad
patterns are more defensible: PyTorch is slowest and most accurate in the aggregate; INT8 is much
faster with a clear accuracy penalty; and the NCNN formats retain most FP32 accuracy while
providing intermediate throughput.

\section{Resource, Counting and Thermal Behaviour}

Peak process memory varied less dramatically than latency. TFLite INT8 had the lowest measured
peak at 824.22~MB. PyTorch, NCNN FP32 and NCNN FP16 were clustered near 835--839~MB, while
TFLite FP32 reached 868.27~MB. The absence of a large memory reduction for every compact numeric
format reinforces that process RSS includes runtime libraries, buffers and interpreter overhead in
addition to model parameters.


Counting behaviour provides an application-facing check on these detection differences. The three
FP32 implementations each produced count MAE 0.600, NCNN FP16 0.611 and TFLite INT8 0.633.
The small spread suggests broad preservation of total-rock counting, but it does not imply that
individual errors are interchangeable. A missed half-buried rock and an additional false-positive
exposed rock can partially cancel in image-level count while still representing two detection
mistakes. Qualitative review of dense and ambiguous scenes is therefore useful when diagnosing
why count error changes, although no unverified example montage is introduced into the main body.

CPU utilisation shows a similarly runtime-specific profile. PyTorch FP32 averaged 44.47\%,
TFLite FP32 69.95\%, TFLite INT8 62.83\%, NCNN FP32 83.32\% and NCNN FP16 87.08\%. A higher
CPU percentage is not automatically worse if it accompanies much lower latency; rather, it
indicates how aggressively the backend uses available processing capacity. The NCNN formats
achieved strong throughput while showing the highest mean CPU utilisation, whereas INT8 achieved
the best throughput without the highest CPU figure.


This matters for system integration. A runtime that occupies most CPU capacity may leave less
headroom for camera handling, GPS processing, communication or control logic even if its raw
inference latency is attractive. Conversely, unused CPU in the PyTorch result does not make it
preferable when the achieved throughput is too low. A future end-to-end system should therefore
measure total application utilisation after all concurrent components are integrated, rather than
assuming the isolated detector profile will remain unchanged.

During the fixed 90-image benchmark, no recorded thermal throttling occurred. Approximate maximum
temperatures were 55.10~$^\circ$C for TFLite INT8, 56.75~$^\circ$C for NCNN FP32,
59.50~$^\circ$C for NCNN FP16, 61.70~$^\circ$C for PyTorch and 66.10~$^\circ$C for TFLite FP32.
These short-run maxima should not be interpreted as equilibrium temperatures.

\reportchart[0.82\textwidth]{figures/figure_11.png}{Mean CPU usage, peak process memory and maximum short-run temperature for each completed format. Faster inference does not translate uniformly into lower CPU load, memory demand or thermal output.}{fig:resource-thermal-chart}

The continuous test rose from the mid-40s to approximately 76--77~$^\circ$C while sustaining
about 22.7 FPS over roughly 2,533 recorded inferences. No throttling or reboot was recorded, but
the upward thermal trajectory means the run is evidence of short-duration stability only. The
Raspberry Pi 5 is designed for substantially higher performance than earlier boards, and official
documentation explicitly treats cooling as important for sustained workloads
\cite{raspberrypi5_docs,raspberrypi_cooling}. A controlled active-cooling experiment is therefore a
clear next step rather than a basis for claiming that the absence of a cooler caused any specific
failure observed here.


The contrast between the 90-image temperatures and the continuous INT8 trace also illustrates why
short benchmarks cannot establish thermal sustainability. INT8 finished the fixed workload near
55~$^\circ$C, yet repeated inference continued to accumulate heat until the upper-70s within about
two minutes. A format that finishes a short test at a moderate temperature could therefore behave
very differently after several minutes. Long-run testing should record CPU frequency and throttle
flags together with temperature so that performance degradation can be distinguished from a
simple rise in sensor temperature.

\section{Operational Stability and Runtime Failures}

The negative runtime results are analytically important because a deployment format that cannot
complete the intended workload cannot be ranked only from isolated inference speed. ONNX FP32
successfully processed a single image, proving that export and basic execution were possible, but
sustained attempts ended in abrupt reboots. No valid 90-image aggregate ONNX benchmark therefore
exists. Potential explanations include thermal, power, memory or runtime behaviour, but the logs
do not isolate a cause; assigning one would exceed the evidence.

NCNN FP32 presents a different lesson. It completed the fixed 90-image benchmark and therefore
has valid aggregate metrics, yet a later continuous test hard-reset the Pi. Short benchmark
completion is consequently weaker evidence than repeated sustained operation. TFLite INT8 is the
only format for which both a complete 90-image benchmark and the genuine approximately two-minute
continuous experiment are available.


The evidence hierarchy can therefore be stated explicitly. A successful export establishes file
conversion; a successful single-image smoke test establishes basic functional compatibility; a
complete 90-image benchmark establishes short-workload aggregate behaviour; and a continuous test
adds evidence about sustained operation. These levels should not be collapsed. ONNX reached the
single-image level, NCNN FP32 reached the fixed-benchmark level but failed a later continuous
attempt, and INT8 reached the available continuous-test level. NCNN FP16 has strong fixed-benchmark
metrics but was not subjected to equivalent continuous validation.

These outcomes support the general benchmarking principle that representative workloads and
system-level measurements matter, rather than assuming portable model formats behave equivalently
on a particular device \cite{reddi2020}. They also strengthen the decision to report failures as
results rather than substitute planned numbers or infer missing aggregate values.


\section{Deployment Selection}

For the tested Raspberry Pi 5 CPU-only configuration, TFLite INT8 provides the strongest combined
evidence. It has the lowest mean latency (44.13~ms), highest median throughput (22.62 FPS), lowest
peak process memory among the completed formats (824.22~MB), a successful 90-image run and the
only completed continuous-inference test. The selection is therefore based on a combination of
speed, resource profile and available stability evidence, not on speed alone.

The recommendation is explicitly contextual. TFLite INT8 has lower accuracy than the FP32 and
FP16 alternatives, and a different hardware accelerator, cooling configuration, workload or
accuracy requirement could alter the preferred point. NCNN FP16 is the strongest secondary
candidate where accuracy retention is more important: it reached 13.82 FPS with mAP50 0.537,
close to the PyTorch mAP50 of 0.548. Its limitation in this study is the absence of an equivalent
completed continuous stability experiment.


A useful way to interpret the final choice is as a Pareto-style decision. PyTorch occupies the
high-accuracy/low-speed end; INT8 occupies the high-speed/lower-accuracy end; and NCNN FP16 lies
between them. No single representation dominates the others on every criterion. The chosen point
reflects the project's edge-computing emphasis and the evidence actually completed on the device.
If the downstream system assigns a very high cost to missed rocks, the decision boundary could
shift toward NCNN FP16 even before its throughput becomes the maximum.


The selection can be stress-tested against alternative priorities. If mAP50 retention were treated
as the dominant objective, NCNN FP16 would be attractive because it loses only 0.011 absolute
mAP50 from PyTorch while increasing throughput by more than fivefold. If maximum throughput and
available continuous evidence dominate, INT8 is stronger. If simplicity and fidelity to the
original training framework dominate, PyTorch avoids conversion but fails the speed objective.
This exercise demonstrates that the recommendation is not a hidden weighted average chosen after
seeing the results; it is a transparent prioritisation of edge responsiveness with an explicit
accuracy constraint.

A production decision would require those priorities to be stated by the agricultural operator.
For example, the cost of missing a rock may depend on machinery damage risk, while a false positive
may cause an unnecessary stop or removal attempt. Such costs are not available in the current
project, so the dissertation avoids converting mAP and count MAE into an unsupported economic
utility score. The recommendation remains technical and configuration-specific.


The same caution applies to the apparent 22.62 FPS headline. Inference throughput is only one
component of an eventual field pipeline. A production system would also need to decode camera
frames, maintain timestamps, transform or associate coordinates, log detections and potentially
communicate with a control system. Those operations consume CPU and memory that were not present
in the stored-image benchmark. The measured value should therefore be read as available detector
throughput under the test workload, not as a guaranteed end-to-end frame rate after all future
components are integrated.

\section{Threats to Validity and Relation to Existing Work}

Several limitations bound the strength of the conclusions. First, the evaluation used a single
Raspberry Pi 5 rather than repeated devices, so manufacturing variation and device-to-device
thermal behaviour are not represented. Second, the executed benchmark used three warm-ups and
one timed pass, not repeated independent timing runs. Tail latency and variance can therefore be
reported from the pass but not as a robust distribution across repeated experiments.

Third, the 90-image test set was untouched for \emph{architecture} selection, but was later reused
to compare and choose the final edge runtime. There is no second untouched dataset that provides
an unbiased final estimate after runtime selection. Fourth, the workload uses stored letterboxed
images rather than a live camera pipeline. Image capture, decoding, frame buffering and field
motion could change end-to-end latency. Fifth, the stability evidence is asymmetric: INT8 has a
completed approximately two-minute continuous run, NCNN FP32 has a failed continuous attempt,
and the other formats do not have equivalent long repeated tests.

Sixth, accuracy results depend on the selected confidence and IoU settings. Alternative operating
points could change precision/recall and count error, especially for half-buried rocks. Finally,
the project does not validate GPS mapping, size/weight estimation or robotic removal. It should be
positioned alongside precision-agriculture edge research as a focused deployment experiment, not
a field-ready autonomous system. This framing is consistent with reviews that identify Raspberry
Pi as a useful low-cost agricultural platform while calling for broader domain and real-world
validation \cite{joice2025}, and with edge-agriculture literature that emphasises unresolved
accuracy/resource trade-offs \cite{sapna2026}.


There is also a distinction between internal and external validity. Internal comparability is
strengthened by fixing the images, detector and thresholds across completed runtimes. External
validity is weaker because the results come from one board, one dataset domain and one stored-image
pipeline. The study is therefore strongest when answering ``which representation behaved best in
this controlled configuration?'' and weaker when answering ``which representation will be best on
all agricultural edge systems?'' The latter claim is intentionally not made.


The choice of test set creates one further interpretive boundary. The 90 images were legitimately
untouched when the detector architecture was selected, but after those same images were used to compare
runtime representations and choose TFLite INT8, the final runtime decision became conditioned on them.
Accordingly, the reported INT8 metrics describe performance on the selection workload, not a fresh
post-selection estimate. A rigorous next experiment would freeze the complete detector-plus-runtime
pipeline and evaluate it once on a second independently collected set. That step would test whether the
measured INT8 compromise and the half-buried weakness persist outside the data used for the final edge
choice.


The findings nevertheless align with several themes in related work. Integer quantisation has long
been motivated as a way to enable efficient inference on resource-constrained processors
\cite{jacob2018}; the measured INT8 result demonstrates that the expected speed benefit can be
substantial on the target Pi, while also showing a non-zero accuracy cost. MLPerf's emphasis on
well-defined inference scenarios is reflected in the discovery that the same YOLOv8n detector
spans roughly 2.5 to 22.6 FPS depending on representation \cite{reddi2020}. In agricultural
context, the persistence of a half-buried weakness is compatible with broader reviews that identify
scene variability and dataset coverage as central challenges \cite{ariza2024,badgujar2024}.

The project also complements work using alternative sensing. Th\"urkow et al. investigate thermal
imagery for agricultural stone detection \cite{thurkow2025}, suggesting a potential route for
cases in which RGB contrast is poor. The present study does not establish that thermal or depth
sensing would solve the burial-status problem, but it provides a concrete reason to test such
modalities: the dominant remaining weakness is semantic visibility of partially buried stones,
not simply insufficient CPU throughput.

\section{Overall Answer to the Research Question}

The two-class YOLOv8n detector can be deployed effectively on a Raspberry Pi 5 for stored-image,
CPU-only inference, but effectiveness depends strongly on representation. PyTorch FP32 provides
the best measured aggregate detection quality among the completed edge runs but is slow at
397.48~ms mean latency and 2.52 median FPS. TFLite INT8 reduces mean latency to 44.13~ms and
raises throughput to 22.62 FPS, while mAP50 decreases from 0.548 to 0.503 and mAP50--95 from
0.275 to 0.234. NCNN FP16 occupies an intermediate position at 13.82 FPS and mAP50 0.537.

The most important accuracy weakness is not solely a conversion artefact: half-buried rocks were
already much harder than completely exposed rocks in the source detector. Operationally, five
representations completed the fixed benchmark, whereas ONNX did not produce a valid aggregate
run. TFLite INT8 also completed the available short continuous experiment without recorded
throttling or reboot, although temperature rose to approximately 76--77~$^\circ$C. The research
question is therefore answered by a trade-off rather than a universal winner: TFLite INT8 is the
preferred format for the tested proof-of-concept configuration, while NCNN FP16 remains a credible
accuracy-retention alternative requiring stronger stability validation.

% =============================================================================
% CHAPTER 5
% =============================================================================

\chapter{Conclusion and Future Work}
\label{chap:conclusion}

\section{Overall Conclusion and Answer to the Research Question}

This dissertation investigated how effectively a two-class YOLOv8n rock burial-status detector
could be transferred from model development to Raspberry Pi 5 CPU inference. The project
established a common 90-image edge workload, prepared six runtime representations and obtained
complete aggregate results for five of them. The evidence shows that model representation has a
substantial effect on practical deployment even when the trained detector is held constant.

PyTorch FP32 produced the strongest measured aggregate edge accuracy, with mAP50 0.548 and
mAP50--95 0.275, but required 397.48~ms mean latency and achieved only 2.52 median FPS. TFLite
INT8 reduced latency to 44.13~ms and increased throughput to 22.62 FPS---approximately a
9.01$\times$ latency speed-up---at the cost of reducing mAP50 to 0.503 and mAP50--95 to 0.234.
NCNN FP16 offered a useful middle point with 13.82 FPS and mAP50 0.537. ONNX FP32 could process
isolated images but did not complete a valid sustained aggregate benchmark.

For the tested CPU-only Raspberry Pi 5 configuration, TFLite INT8 is therefore the preferred
representation because it combines the best measured latency/throughput, the lowest peak process
memory among completed formats, a complete 90-image benchmark and the available short continuous
stability evidence. This is not a claim that INT8 is universally superior. The appropriate runtime
would change if accuracy requirements, hardware acceleration, cooling or field conditions changed.

\section{Assessment of Project Objectives and Research Hypothesis}

\begin{table}[htbp]
\centering
\caption{Assessment of the project objectives against completed evidence.}
\label{tab:objectives}
\begin{tabular}{|>{\raggedright\arraybackslash}p{8.8cm}|>{\raggedright\arraybackslash}p{5.2cm}|}
\hline
Objective & Status \\
\hline
Establish the two-class detector and evaluation context & Achieved \\
\hline
Prepare six alternative edge-runtime representations & Achieved \\
\hline
Create a reproducible Raspberry Pi evaluation workflow & Achieved \\
\hline
Benchmark quality, latency, throughput and resource behaviour & Achieved for five formats; ONNX aggregate incomplete \\
\hline
Investigate operational reliability and negative results & Partially achieved; short/asymmetric stability evidence \\
\hline
Select and justify a preferred representation for the tested Pi & Achieved: TFLite INT8 \\
\hline
\end{tabular}
\end{table}

The research hypothesis is supported. Every completed deployment-oriented representation reduced
mean latency relative to PyTorch FP32, with speed-ups ranging from approximately 2.77$\times$ for
TFLite FP32 to 9.01$\times$ for TFLite INT8. The magnitude of improvement varied substantially by
runtime. The hypothesis also anticipated possible accuracy or robustness costs: INT8 showed the
largest accuracy reduction, ONNX failed to complete sustained benchmarking, and NCNN FP32 failed
a later continuous attempt despite completing the fixed benchmark. The evidence therefore
supports the proposed accuracy--efficiency--robustness trade-off rather than a simple claim that
lower precision is always better.

\section{Current Project Status and Limitations}

The resulting system is an academic edge-deployment proof of concept. It demonstrates stored-image
rock detection and burial-status classification on a Raspberry Pi 5, not a complete autonomous
field solution. Live-camera acquisition, GPS integration, verified rock size/weight estimation,
3D sensing and robotic removal remain outside the implemented scope described in the wider project
materials \cite{ccrah_rock_picking,campelo2026kickoff}.

The strongest model-level limitation is half-buried recognition. The source detector already
showed a large gap between completely exposed and half-buried rocks, so this issue should not be
attributed solely to INT8 or FP16 conversion. Experimental validity is also constrained by the
single Raspberry Pi, one timed pass per completed format, fixed thresholds, stored-image workload,
short stability evidence and reuse of the 90-image set for final runtime selection after it had
served as the untouched architecture-selection test. A second untouched deployment-validation set
would be required for a cleaner final estimate of the selected runtime.

\section{Prioritised Future Work}

The first priority is to improve half-buried detection. Additional independently collected examples
should cover different soil types, lighting, moisture, occlusion levels, rock textures and dense
scenes. Hard-negative mining and targeted augmentation should be evaluated against a held-out field
set rather than inferred from training behaviour alone.

Second, the runtime comparison should be repeated with multiple timed passes and longer stability
experiments under controlled ambient conditions. Active cooling should be evaluated directly,
with temperature, frequency/throttling state, power and latency recorded throughout. NCNN FP16
should receive the same continuous protocol as TFLite INT8 because its accuracy retention makes it
a strong alternative.

Third, ONNX failure diagnosis should be made systematic. Repeated runs should isolate memory,
power, thermal and runtime variables rather than inferring a cause from an abrupt reboot. Only
a completed controlled run should be used to add ONNX to the aggregate ranking.

Fourth, evaluation should progress from stored images to a live-camera pipeline. End-to-end
measurements should include frame acquisition, resizing/letterboxing, inference, post-processing
and output handling. Once the detector is stable in that setting, the wider project can consider
GPS association, rock-size estimation, multimodal or depth sensing, and eventual integration with
robotic removal. These stages should be treated as new validation problems rather than assumed to
follow automatically from successful detector inference.

\section{Final Conclusion}

The project demonstrates that a lightweight agricultural detector can run at useful speed on a
Raspberry Pi 5, but deployment decisions materially change accuracy, latency, resource demand and
operational confidence. TFLite INT8 provides the strongest overall evidence for the tested
CPU-only proof of concept, while NCNN FP16 preserves more accuracy at lower throughput. The most
important next challenge is not simply to make inference faster: it is to strengthen half-buried
rock recognition and validate sustained, cooled, live-camera operation under realistic field
conditions. Collectively, the contributions described here are outcomes of the group project;
individual roles, decisions and learning are documented separately in the Individual Reflective
Accounts.

% =============================================================================
% BACK MATTER
% =============================================================================

\clearpage

% Preserve the University-provided BibTeX workflow.
% Replace/extend the entries inside sample_bibtex.bib with the final verified
% project bibliography while keeping the database call below unless deliberately
% renamed later.
\bibliography{sample_bibtex.bib}

% -----------------------------------------------------------------------------
% APPENDICES
% -----------------------------------------------------------------------------

\appendix

\chapter{Dataset and Detector Development Evidence}
\label{appx:detector}

The original dataset contained 448 images and 706 labelled rocks: 362 completely exposed and
344 half buried. The pre-augmentation training/validation/test split is reported in
Table~\ref{tab:dataset}. Training-only augmentation generated 1,878 additional images, producing
2,191 training images and 3,563 annotation instances across the augmented copies.

\begin{table}[htbp]
\centering
\caption{Extended detector-development evidence.}
\begin{tabular}{|l|r|r|r|r|r|}
\hline
Model & Precision & Recall & F1 & mAP50 & mAP50--95 \\
\hline
YOLOv8n & 0.681 & 0.759 & 0.718 & 0.759 & 0.4288 \\
\hline
YOLO11n & 0.743 & 0.719 & 0.731 & 0.729 & 0.4002 \\
\hline
YOLO11s & 0.764 & 0.801 & 0.782 & 0.783 & 0.4107 \\
\hline
\end{tabular}
\end{table}

The selected YOLOv8n produced untouched-test mAP50--95 of approximately 0.3384. At the evaluated
source-detector operating point, the class-specific confusion result represented 64/69 completely
exposed rocks correctly and 23/66 half-buried rocks correctly. These values are operating-point
correct rates and must not be relabelled as AP.

\chapter{Raspberry Pi Deployment and Reproducibility}
\label{appx:deployment}

The six prepared representations were PyTorch FP32, ONNX FP32, TFLite FP32, TFLite INT8,
NCNN FP32 and NCNN FP16. Final inference used fixed $640\times640$ letterboxed images with padding
value 114, batch size one, confidence threshold 0.25 and IoU setting 0.70. Three warm-up inferences
preceded one timed 90-image pass. The ten-image smoke test was a subset of the same 90 images and
served only as a functional check.

INT8 calibration used 200 original training images containing 304 labelled objects, balanced at
152 completely exposed and 152 half buried. The Raspberry Pi environment was Debian GNU/Linux 13,
Python 3.13.5, Ultralytics 8.4.115, ONNX Runtime 1.28.0, OpenCV 5.0.0, NumPy 2.5.1,
Pandas 3.0.5 and psutil 7.2.2. The device used the official 27~W 5.1~V/5~A supply and no Active
Cooler. No camera was used in the benchmark.

\chapter{Extended Raspberry Pi Results}
\label{appx:results}

\begin{table}[htbp]
\centering
\caption{Extended canonical edge metrics used in the main evaluation.}
\resizebox{\textwidth}{!}{%
\begin{tabular}{|l|r|r|r|r|r|r|r|r|r|r|}
\hline
Format & Latency ms & FPS & CPU \% & RAM MB & P & R & F1 & mAP50 & mAP50--95 & Count MAE \\
\hline
PyTorch FP32 & 397.48 & 2.52 & 44.47 & 839.47 & 0.716 & 0.550 & 0.622 & 0.548 & 0.275 & 0.600 \\
\hline
TFLite FP32 & 143.70 & 6.95 & 69.95 & 868.27 & 0.669 & 0.542 & 0.599 & 0.535 & 0.270 & 0.600 \\
\hline
TFLite INT8 & 44.13 & 22.62 & 62.83 & 824.22 & 0.661 & 0.539 & 0.594 & 0.503 & 0.234 & 0.633 \\
\hline
NCNN FP32 & 75.09 & 13.26 & 83.32 & 835.25 & 0.667 & 0.542 & 0.598 & 0.535 & 0.271 & 0.600 \\
\hline
NCNN FP16 & 72.31 & 13.82 & 87.08 & 835.45 & 0.672 & 0.542 & 0.600 & 0.537 & 0.270 & 0.611 \\
\hline
ONNX FP32 & N/A & N/A & N/A & N/A & N/A & N/A & N/A & N/A & N/A & N/A \\
\hline
\end{tabular}}
\end{table}

Approximate maximum temperatures during the fixed benchmark were 55.10, 56.75, 59.50, 61.70 and
66.10~$^\circ$C for TFLite INT8, NCNN FP32, NCNN FP16, PyTorch FP32 and TFLite FP32 respectively.
These values are short-workload observations and not long-run thermal equilibria.

\chapter{Stability and Failure Diagnostics}
\label{appx:stability}

TFLite INT8 completed the only valid continuous-inference experiment retained for final analysis:
approximately two minutes, roughly 2,533 inference rows and about 22.7 FPS. Temperature rose from
the mid-40s to about 76--77~$^\circ$C (raw maximum approximately 77.1~$^\circ$C), with no recorded
reboot or throttling during the run.

ONNX FP32 processed isolated images but did not complete a valid 90-image aggregate benchmark;
sustained attempts caused abrupt Raspberry Pi reboots and the root cause remains unresolved. NCNN
FP32 completed the 90-image benchmark but a later continuous attempt hard-reset the device and did
not yield a valid completed continuous CSV. No genuine 10-minute or 30-minute final stability test
was completed. A file whose name suggested a 30-minute test duplicated the genuine two-minute data
and is excluded from the evidence base.

% =============================================================================

\end{document}
